% Source: https://artofproblemsolving.com/community/c3h2122556_counting_problem_i_made
% Stable identity: aops-topic-2122556-q01
Think of each of the first six customers as choosing a type independently, each type with probability $\dfrac{1}{3}$. As long as fewer than three refrigerators have been sold, a customer who chooses a refrigerator buys one. So all three are sold within the first six purchases exactly when at least $3$ of the first $6$ customers choose a refrigerator.

The number of customers among the six who choose a refrigerator is binomial with $n=6$ and $p=\dfrac{1}{3}$, so the probability is
\[\sum_{k=3}^{6}\dbinom6k\left(\dfrac{1}{3}\right)^k\left(\dfrac{2}{3}\right)^{6-k}=\dfrac{\dbinom63\cdot8+\dbinom64\cdot4+\dbinom65\cdot2+\dbinom66}{3^6}=\dfrac{160+60+12+1}{729}=\boxed{\dfrac{233}{729}}\approx0.32 .\]
(Check by the complement: $\dfrac{64+192+240}{729}=\dfrac{496}{729}$ for at most two refrigerator choices, and $729-496=233$.)
